\section{危害公共安全罪 - 危险方法型、破坏型、交通型犯罪}
\begin{enumerate}
  \item 社会法益罪名最多，占 80\%，但分值少。
  \item 危害。危险和实害。产生危险是成立，造成实害是既遂。
  \item 公共。危及范围无法预料，不特定多数人。
  \item 安全。人身伤亡、财产损失，不保护精神、心理安全。
  \item 标准。\zd{危及公共多数人的可能性}。\al{扔砖块、扔手雷、楼顶泼硫酸}。
  \item 危险方法型犯罪。
        \begin{enumerate}
          \item 放火、投放危险物质。
                \begin{enumerate}
                  \item \al{甲想报复单位，认为仓库没人把守，放火烧了仓库，烧死库管员}。
                  \item \al{甲想杀死乙，向乙家放火，导致乙死亡和邻居房屋被毁}。
                  \item \al{甲基于仇富心理，大巴放火烧死乘客}。
                  \item \al{甲放火烧死乙全家，乙独门独户别墅，四周无人}。
                  \item \al{宿舍饮水机投毒；楼道尽头饮水机投毒}。
                \end{enumerate}
          \item 以危险方法危害公共安全。\zd{不是兜底罪名，是补充罪名，最后适用}。
                \begin{enumerate}
                  \item 具体的、现实紧迫的危险。\al{甲盗窃消防栓铜芯}。
                  \item \al{甲寄面粉给乙，谎称是炭疽热病菌粉，乙吓得不轻}。
                  \item 危险的样态。\al{政府大楼每层捅 1 个连捅 6 人}。
                  \item 危险的程度。\al{高速公路点小火堆；高速公路告诉逆行；盗窃井盖}。
                \end{enumerate}
        \end{enumerate}
  \item 破坏型犯罪。
        \begin{enumerate}
          \item 破坏交通工具。
                \begin{enumerate}
                  \item 交通工具范围。人力三轮、马车、摩托车这些小，破坏了\zd{危害程度没有那么大}，所以不包括。
                  \item 破坏程度要求。\al{砸公交车玻璃}，不影响驾驶安全；\al{拆轮子}，已无法上路；\al{飞机还在停机坪砸玻璃}。
                \end{enumerate}
          \item 破坏交通设施。物理 + 效用/功能用途丧失。
        \end{enumerate}
  \item 交通型犯罪。最重要的。
        \begin{enumerate}
          \item 交通肇事。
                \begin{enumerate}
                  \item 主体\zd{不限于司机}。\al{乙喝醉跟甲借车出事}。
                  \item 地点要求公共交通领域。\al{建筑工地/大学校园内把人撞死了}。
                  \item 行为必须是过失。
                  \item 结果。\yi 实害结果（死 1 \hz 重 3 \hz 重 1 + 严重情节（醉驾、毒驾）\er 因果关系。
                  \item 法定刑升格条件。
                        \begin{enumerate}
                          \item 一般 \yjt\, 3 年以下。
                          \item 肇事逃逸 \yjt\, 3 - 7。\zd{要意识到出事了}。\al{甲喝得不省人事撞死乙，拖着丙开了很远，被人拦住才知}。
                          \item 因逃逸致人死亡 \yjt\, 7 - 15。
                                \begin{itemize}
                                  \item \zd{本质是不作为犯、结合犯}。交通肇事 + 遗弃罪致人死亡 = 交通肇事加重处罚。
                                  \item \al{甲撞伤乙，乙亲戚群涌而上围殴，甲只能驾车跑到公安局报案}。
                                  \item 逃逸和死亡要有因果关系。肇事行为 \yjt 不救人不作为 \yjt 介入因素 \yjt 死亡结果。
                                  \item \al{连环碾压案}。无法查明，死得越早对“大家”越有利。
                                \end{itemize}
                        \end{enumerate}
                  \item 共犯。
                        \begin{enumerate}
                          \item 共同过失。\al{乙喝醉跟甲借车出事}。甲交通肇事罪的帮助犯 ×。
                          \item \al{指使逃逸}。\al{甲撞伤乙，丙说救啥救啊，赶紧跑吧}。
                                \begin{itemize}
                                  \item 丁 \yjt 遗弃罪教唆犯、窝藏罪（扩大解释，帮他跑路，不是藏在家里才是窝藏）。
                                  \item 指使逃逸的人是\zd{单位主管人员、车辆所有人、承包人、乘车人}定交通肇事罪共犯。
                                \end{itemize}
                        \end{enumerate}
                \end{enumerate}
          \item 危险驾驶。
                \begin{enumerate}
                  \item \zd{醉酒驾驶}和酒后驾驶。血液含量 100/80。\zd{不要求主观上认识到，概括认识}。
                  \item \al{喝酒突发心脏病开车送医院} \yjt 紧急避险。
                  \item 共犯问题。
                        \begin{enumerate}
                          \item \al{乙喝醉跟甲借车出事}。甲危险驾驶罪的帮助犯 √。
                          \item \al{甲准备出狱了，让乙来看望他，乙醉酒开车去了}。
                          \item \al{甲喝醉酒还去运输毒品}。
                        \end{enumerate}
                \end{enumerate}
          \item 妨害安全驾驶。\zd{行驶中}；公共交通工具；打司机 \hz 抢方向盘。
        \end{enumerate}
\end{enumerate}
